\section{Appendix: Identification and observable risk certificates}

The auxiliary results connect the structural causal model to population localization and observable estimation. The argument begins with conditional means at the realized dose, then establishes uniqueness under Gaussian convolution and positivity of localized population denominators. A public expected-error certificate subsequently links these population quantities to the split-sample estimator, while dictionary comparisons supply the uniform rate guarantees in \cref{obj:thm:observable-upper,obj:thm:finite-honesty}.

\subsection{Random-dose means and identification}

Conditional unconfoundedness in \cref{obj:ass:conditional-unconfoundedness} concerns the measurable potential-outcome schedule. The regression identity needed for identification and weighting evaluates that schedule at the random latent dose. The following result supplies this bridge, including its integrated form for bounded measurable multipliers.

\begin{lemmav}{lem:realized-dose-mean}[Realized dose mean identities]
Suppose the following conditions hold:
\begin{itemize}
\item \textbf{(Public model primitives.)} Fix a potential-outcome path space as in \cref{obj:def:potential-path-space} and public class constants \(K\) as in \cref{obj:def:model-class}.
\item \textbf{(Parameter ranges.)} Let \(\beta\in(0,1]\), \(\kappa\in[0,2]\), and \(\sigma\in[0,1/4]\).
\item \textbf{(Structural law.)} Let \(P\) be a structural probability law in the model class \(\mathcal P_{K,\beta,\kappa,\sigma}\) of \cref{obj:def:model-class}. Write \(\mathcal X=\{0,1\}\) for the stratum space, and let \(X\in\mathcal X\), \(A\in\mathbb R\), \(Z\in\mathbb R\), and \(Y\in\mathbb R\) denote its stratum, latent dose, classical error, and realized outcome coordinates, respectively.
\end{itemize}
With \(\mu_x(a)\) denoting the conditional potential-outcome mean defined in \cref{obj:def:conditional-potential-mean}, the following statements hold \(P\)-almost surely:
\[
0\le Y\le1,\qquad
E_P[Y\mid A,X,Z]=\mu_X(A),\qquad
E_P[Y\mid A,X]=\mu_X(A).
\]
Moreover, for every bounded measurable function
\(F:\mathbb R\times\mathcal X\times\mathbb R\to\mathbb R\),
\[
E_P\!\left[F(A,X,Z)Y\right]
=
E_P\!\left[F(A,X,Z)\mu_X(A)\right].
\]
\end{lemmav}

The conditional identities make the potential-outcome mean applicable to the observed outcome at the realized dose, including conditioning on the error coordinate. The almost-sure outcome bound also supplies the boundedness used in the observable moment certificates below. Joint measurability in \cref{obj:def:potential-path-space} gives random-dose evaluation its measurable meaning.

For the convolution argument, let \(\Phi_\sigma\) denote the centered Gaussian probability measure with variance \(\sigma^2\). At zero error scale, set \(\Phi_0=\delta_0\), where \(\delta_0\) is the unit point mass at zero. For a finite Borel measure \(\lambda\), its convolution with this error measure is defined, for each Borel set \(B\subseteq\mathbb R\), by
\[
(\lambda*\Phi_\sigma)(B)
=
\int_{\mathbb R}\int_{\mathbb R}
\mathbf 1\{t+u\in B\}\,d\Phi_\sigma(u)\,d\lambda(t).
\]
The centered latent-dose support in \cref{obj:ass:latent-support} motivates the following uniqueness statement for finite measures supported on \([-1/2,1/2]\).

\begin{lemmav}{lem:compact-gaussian-convolution-injective}[Injectivity of compact Gaussian convolution]
Let \(\sigma\in\mathbb R\), and let \(\Phi_\sigma\) denote the centered Gaussian probability measure with variance \(\sigma^2\), with \(\Phi_0=\delta_0\), the unit point mass at zero. Let \(\mu_1,\mu_2,\nu_1,\nu_2\) be nonnegative Borel measures on \(\mathbb R\). Suppose that:
\begin{itemize}
\item \textbf{(Error scale.)} \(0\le \sigma\le 1/4\).
\item \textbf{(Finiteness and support.)} Each of \(\mu_1,\mu_2,\nu_1,\nu_2\) is finite and assigns zero mass to \(\mathbb R\setminus[-1/2,1/2]\).
\item \textbf{(Convolution equality.)}
\[
\mu_1*\Phi_\sigma+\nu_2*\Phi_\sigma
=
\nu_1*\Phi_\sigma+\mu_2*\Phi_\sigma.
\]
\end{itemize}
Then
\[
\mu_1+\nu_2=\nu_1+\mu_2.
\]
\end{lemmav}

The grouped equality accommodates differences of finite nonnegative measures while retaining a statement entirely in terms of such measures. Its mathematical role is uniqueness of latent measures from their Gaussian-convolved counterparts. The formulation includes the error-free case through the point-mass convention.

Population localization additionally requires a positive denominator. Here positivity belongs to the latent weight \(q\), whereas empirical moments use its observable inverse \(\ell_q\). Exact Gaussian inversion connects the two under the stated integrability conditions. The next result records the denominator and approximation bounds needed for this construction.

\begin{lemmav}{lem:positive-ratio}[Positive denominators and approximation]
Fix a potential-path space \((\mathcal S,\mathcal E)\) as in \cref{obj:def:potential-path-space}. Let \(\mathcal X\) consist of two strata, and let \(P\) be a structural probability law with coordinates \(X\in\mathcal X\), \(A,Z,Y,U\in\mathbb R\), and potential-outcome schedule \(Y(\cdot)\in\mathcal S\). Set
\[
a_0=\frac12,\qquad t=A-a_0,\qquad W=A+\sigma Z,\qquad V=W-a_0=t+\sigma Z.
\]
Let \(q:\mathbb R\to\mathbb R\) be a public latent weight and \(\ell_q:\mathbb R\to\mathbb R\) its observable inverse weight. Define the weighted moments by
\[
I_\gamma(q)=\int_{-1/2}^{1/2}|t|^\gamma q(t)\,dt.
\]
Suppose the following conditions hold:
\begin{itemize}
\item \textbf{(Public constants.)} The constants \(K=(p_{\min},c_{\mathrm{lo}},c_{\mathrm{hi}})\) satisfy
\[
0<p_{\min}\le\frac12,\qquad 0<c_{\mathrm{lo}}\le c_{\mathrm{hi}}.
\]
\item \textbf{(Parameter ranges.)}
\[
\beta\in(0,1],\qquad \kappa\in[0,2],\qquad \sigma\in[0,1/4].
\]
\item \textbf{(Structural model.)} The law \(P\) satisfies \cref{obj:def:model-class} with path space \((\mathcal S,\mathcal E)\), constants \(K\), and parameters \((\beta,\kappa,\sigma)\).
\item \textbf{(Admissible weight.)} The functions \(q\) and \(\ell_q\) are measurable, and
\[
q(t)\ge0\quad\text{for every }t\in[-1/2,1/2],
\qquad 0<I_\kappa(q)<\infty.
\]
\item \textbf{(Exact Gaussian inversion.)} For every \(t\in[-1/2,1/2]\), the function \(z\mapsto\ell_q(t+\sigma z)\) is absolutely integrable under the standard Gaussian law, and
\[
\int_{\mathbb R}\ell_q(t+\sigma z)\frac{e^{-z^2/2}}{\sqrt{2\pi}}\,dz=q(t).
\]
The inverse weight is also absolutely integrable under the structural law:
\[
E_P[|\ell_q(V)|]<\infty.
\]
\end{itemize}
For every \(x\in\mathcal X\), the denominator is strictly positive and satisfies
\[
0<E_P[\mathbf1\{X=x\}\ell_q(V)],
\qquad
E_P[\mathbf1\{X=x\}\ell_q(V)]
\ge p_{\min}c_{\mathrm{lo}}I_\kappa(q).
\]
Consequently, the localized population ratio
\[
m_{x,q}
=
\frac{E_P[\mathbf1\{X=x\}Y\ell_q(V)]}
     {E_P[\mathbf1\{X=x\}\ell_q(V)]}
\]
is well defined and, with \(\mu_x(a_0)\) as in \cref{obj:def:conditional-potential-mean}, obeys
\[
|m_{x,q}-\mu_x(a_0)|
\le
\frac{c_{\mathrm{hi}}}{c_{\mathrm{lo}}}
\frac{I_{\kappa+\beta}(q)}{I_\kappa(q)}.
\]
These conclusions hold throughout the stated parameter ranges, including \(\sigma=0\).
\end{lemmav}

The lower bound \(b_q\) makes population normalization explicit, and \(B_q\) quantifies localization around the target dose. Both certificates use moments of the nonnegative latent weight. Observable inverse weights are governed instead by inversion and integrability; their empirical denominators receive the public floor specified in \cref{obj:synth_3}. Thus population positivity and empirical stabilization have distinct roles.

The realized-dose identity and convolution injectivity support the identification conclusion in \cref{obj:prop:identification}: the first connects outcome-weighted latent measures to conditional potential means, and the second supplies uniqueness under the known measurement channel. The shrinking-neighborhood ratio statement of that proposition is proved directly in steps 3 and 4 of its proof. \Cref{obj:lem:positive-ratio} serves the estimation analysis instead: its weighted-denominator lower bound and localization-bias bound are the population inputs to the error certificate in \cref{obj:lem:public-error-certificate}.

\subsection{Split-sample error certification}

The observable construction in \cref{obj:synth_3} assigns separate modulo-three blocks to stratum frequencies, inverse-weight denominators, and outcome numerators. Its public denominator floor and projection define the empirical ratios for every sample with nonempty blocks. The following certificate quantifies expected absolute error for a specified weight pair before dictionary selection.

\begin{lemmav}{lem:public-error-certificate}[Public expected-error certificate]
Fix the public potential-path space of \cref{obj:def:potential-path-space} and public class constants \(K\) satisfying the restrictions in \cref{obj:def:model-class}. Suppose that:
\begin{itemize}
\item \textbf{(Parameter ranges.)} \(\beta\in(0,1]\), \(\kappa\in[0,2]\), \(n\ge3\), and \(\sigma\in[0,1/4]\).
\item \textbf{(Structural law.)} The structural law \(P\) belongs to \(\mathcal P_{K,\beta,\kappa,\sigma}\), the model class with constants \(K\) defined in \cref{obj:def:model-class}.
\item \textbf{(Sampling.)} The observed sample satisfies \cref{obj:ass:iid}, with joint law \(Q_P^n\) defined in \cref{obj:def:observed-experiment}.
\item \textbf{(Inverse pair.)} Let \(q\) be a public latent weight and \(\ell_q\) its observable inverse weight. Under the usual measurability and integrability conditions, they satisfy
\[
q(t)\ge0\qquad(t\in[-1/2,1/2]),
\qquad
0<I_\kappa(q):=\int_{-1/2}^{1/2}|t|^\kappa q(t)\,dt<\infty,
\]
and the pointwise Gaussian inversion identity
\[
\int_{\mathbb R}\ell_q(t+\sigma z)
       \frac{e^{-z^2/2}}{\sqrt{2\pi}}\,dz
=q(t)
\qquad(t\in[-1/2,1/2]),
\]
where the Gaussian integrand is absolutely integrable for every such \(t\).
\item \textbf{(Law integrability.)} Let \(W\) denote the contaminated dose and \(a_0\) the prespecified evaluation dose, and define the observed-dose displacement by
\[
V:=W-a_0.
\]
Then
\[
\mathbb E_P[|\ell_q(V)|]<\infty.
\]
\item \textbf{(Public second moment.)} Define the public inverse-weight second moment, as an extended nonnegative integral, by
\[
V_q:=
16\int_{-1/2}^{1/2}|t|^\kappa
   \left(
   \int_{\mathbb R}\ell_q(t+\sigma z)^2
       \frac{e^{-z^2/2}}{\sqrt{2\pi}}\,dz
   \right)dt.
\]
Assume \(V_q<\infty\).
\end{itemize}
Let \(\widehat\theta_q\) be the weight-indexed estimator of \cref{obj:def:total-estimator}, and let \(A_q(K)\) be its public expected-error score from that construction, evaluated at the fixed parameters \(\beta,\kappa,n,\sigma\). Define the population causal target by
\[
\theta(P):=\mathbb E_P[Y(a_0)],
\]
where \(Y(a_0)\) is the potential outcome at the evaluation dose in \cref{obj:def:potential-path-space}. Then
\[
\mathbb E_{Q_P^n}
 \left[|\widehat\theta_q-\theta(P)|\right]
\le A_q(K).
\]
\end{lemmav}

In the inverse-pair clause, the ``usual measurability and integrability conditions'' are the following hypotheses: \(q\) and \(\ell_q\) are Borel measurable, and \(t\mapsto|t|^\kappa q(t)\) is Lebesgue integrable on \([-1/2,1/2]\). Together with the displayed nonnegativity, positivity of \(I_\kappa(q)\), and Gaussian inversion with absolute integrability, they form the admissibility conditions on the pair.

The certificate separates localization error, inverse-weight sampling variability, and stratum-frequency error. In particular, the variability term involves the observable inverse through \(V_q(K)\), while localization uses the latent weight through \(B_q(K)\). The denominator bound \(b_q(K)\) connects population normalization to the stability of the clipped empirical ratio. For a public weight pair, all three terms in \(A_q(K)\) are available for the deterministic comparison rule in \cref{obj:def:total-estimator}.

\subsection{Dictionary admissibility and rate comparisons}

The Fourier and polynomial pairs in \cref{obj:synth_1,obj:synth_2} supply two ways to localize while inverting the Gaussian channel. Their finite dictionary also contains the constant pair, so the public minimization rule has a candidate for every \(n\ge1\); the estimator and interval use the selected entry only for \(n\ge3\), when all three sample blocks are nonempty. The next statement verifies the conditions required by the error certificate for every dictionary entry and records the comparisons that control the selected score.

\begin{lemmav}{lem:dictionary-score-rate}[Dictionary admissibility and score rate]
Fix the following public objects and parameters:
\begin{itemize}
\item \textbf{(Path space.)} A potential-path space \((\mathcal S,\mathcal E)\) as in \cref{obj:def:potential-path-space}.
\item \textbf{(Class constants.)} \(K=(p_{\min},c_{\mathrm{lo}},c_{\mathrm{hi}})\), with
\[
0<p_{\min}\le \tfrac12,
\qquad
0<c_{\mathrm{lo}}\le c_{\mathrm{hi}}.
\]
\item \textbf{(Exponents.)} \(\beta\in(0,1]\) and \(\kappa\in[0,2]\).
\end{itemize}
Use the public dictionary \(\mathcal D_{n,\sigma}\) from \cref{obj:def:weight-dictionary} and the public expected-error scores \(A_q(K)\) from \cref{obj:synth_4}. For every integer \(n\ge1\), every \(\sigma\in[0,1/4]\), and every inverse pair \((q,\ell_q)\in\mathcal D_{n,\sigma}\), both weights are measurable, \(q\) is nonnegative on \([-1/2,1/2]\), and its design-weighted moment satisfies
\[
0<I_\kappa(q)
=
\int_{-1/2}^{1/2}|t|^\kappa q(t)\,dt
<\infty.
\]
For every \(t\in[-1/2,1/2]\), the exact Gaussian inversion identity holds with absolute integrability:
\[
\int_{\mathbb R}|\ell_q(t+\sigma z)|\,\mathcal N(0,1)(dz)<\infty,
\qquad
\int_{\mathbb R}\ell_q(t+\sigma z)\,\mathcal N(0,1)(dz)=q(t),
\]
where \(\mathcal N(0,1)\) denotes the standard Gaussian probability measure. The inverse-weight second-moment integral satisfies
\[
16\int_{-1/2}^{1/2}|t|^\kappa
\left(\int_{\mathbb R}\ell_q(t+\sigma z)^2\,\mathcal N(0,1)(dz)\right)\,dt
<\infty.
\]
Let \(V\) denote the observed-dose displacement from the prespecified evaluation dose \(a_0\), and let \(W\) denote the contaminated dose:
\[
V=W-a_0,
\qquad
W=A+\sigma Z,
\]
where \(A\) is the latent dose and \(Z\) is the standardized Gaussian error. For every structural law \(P\in\mathcal P_{K,\beta,\kappa,\sigma}\) from \cref{obj:def:model-class}, the inverse weight is absolutely integrable under the observed centered dose:
\[
\mathbb E_P[|\ell_q(V)|]<\infty.
\]

Let \(\widehat q\) be the first public-score minimizer in the fixed dictionary order from \cref{obj:def:total-estimator}, and let \(\rho_{n,\sigma}\) and \(L_n\) be the frontier rate and logarithmic sample-size quantity from \cref{obj:def:frontier-rate}. There exist \(C>0\) and \(n_0\in\mathbb N\), uniform in \(\sigma\in[0,1/4]\), such that for every \(n\ge n_0\) and every \(\sigma\in[0,1/4]\),
\[
A_{\widehat q}(K)\le C\rho_{n,\sigma}.
\]
For the same \(C\) and \(n_0\), whenever \(\sigma\le L_n^{-1/2}\), there exists \(k\in\mathbb N\) such that the Fourier inverse pair indexed by \(k\) in \cref{obj:def:weight-dictionary} belongs to \(\mathcal D_{n,\sigma}\) and its public score is at most \(C\rho_{n,\sigma}\). Whenever \(L_n^{-1/2}<\sigma\), there exists \(j\in\mathbb N\) such that the polynomial inverse pair with index
\[
m=2^j+1
\]
belongs to \(\mathcal D_{n,\sigma}\) and satisfies
\[
A_{q_m^{\mathrm M}}(K)\le C\rho_{n,\sigma}.
\]
\end{lemmav}

Admissibility holds for every positive sample size, whereas the score comparison holds beyond a threshold independent of the error scale. The Fourier comparison covers \(\sigma\le L_n^{-1/2}\), and the polynomial comparison covers \(L_n^{-1/2}<\sigma\), within the stated large-sample guarantee. These comparisons concern candidates in the same dictionary, so the selected score reflects the applicable localization family through a single public ordering and minimization rule.

\subsection{Observable achievability and interval calibration}

The certificate and dictionary bounds provide the mathematical components of the observable upper bound in \cref{obj:thm:observable-upper}: the certificate controls the error associated with an admissible entry, and the comparison bounds control the selected public score. That result states the expected absolute-error rate uniformly over the public error-scale range and the declared measurable potential-path spaces, for samples beyond its common threshold. The balance is between latent localization and the sampling variability of observable Gaussian inversion.

For interval construction, the same score determines the radius \(A_{\widehat q}(K)/\alpha\) in \cref{obj:def:honest-interval}. The finite-sample certificate and Markov's inequality supply the error control used for calibration, while intersection with \([0,1]\) incorporates the target range. Empty modulo-three blocks trigger the full target-range interval. Accordingly, \cref{obj:thm:finite-honesty} gives coverage of at least \(1-\alpha\) for every \(\alpha\in(0,1)\), every integer \(n\ge0\), every \(\sigma\in[0,1/4]\), and every structural law in the specified model class. Its expected-length bound applies beyond a common sample-size threshold and has order \(\rho_{n,\sigma}\), with a constant depending on \(\alpha\). Coverage and the uniform large-sample precision guarantee therefore retain their respective quantifiers.

\paragraph{Verification note.}
The theorem statements and proofs are machine-checked in Lean 4. The definition blocks written for exposition (\cref{obj:synth_1,obj:synth_2,obj:synth_3,obj:synth_4,obj:synth_5,obj:synth_6,obj:synth_7,obj:synth_8}) restate in prose the Lean definitions and constructions that the theorems use; these prose blocks are not themselves machine-checked statements.
